\section{Introduction}\label{sec:introduction}

For $\abs{x}\le 1$ let
\[
  \chi_2(x)=\sum_{j\ge0}\frac{x^{2j+1}}{(2j+1)^2}=\tfrac12\bigl(\Li_2(x)-\Li_2(-x)\bigr),
  \qquad
  \Ti_2(x)=\sum_{j\ge0}\frac{(-1)^jx^{2j+1}}{(2j+1)^2}=\operatorname{Im}\Li_2(ix)
\]
be Legendre's chi function and the inverse tangent integral. This paper presents a proof of the
following statement.

\begin{theorem}\label{thm:main}
For every odd integer $q\ge 81$ the numbers $\chi_2(2/q)$ and $\Ti_2(2/q)$ are irrational.
Since $\chi_2$ and $\Ti_2$ are odd, the same holds for $\chi_2(-2/q)$ and $\Ti_2(-2/q)$.
\end{theorem}

\emph{Status.} The proof is a candidate. It was reconstructed and checked by an AI system from
candidate arguments produced by another AI system (see the disclosure at the end); no
independent human expert has reviewed it. In Lean only the final deduction is formalised, and
only conditionally on the three estimates of \cref{thm:three-estimates}. All finite
inequalities used in the proof have been recomputed by second implementations, also written by an
AI system, that do not import the code that produced them (Appendices~\ref{app:certificates}
and~\ref{app:consistency}); this checks the numbers, not the asymptotic arguments, and the two
implementations share inputs and authorship, so it is not an independent verification. \Cref{subsec:step-status} gives the status of each
step, and \cref{rem:scrutiny} lists the points that most need independent scrutiny.

\subsection{Background}\label{subsec:background}

For even $q$ one has $\chi_2(2/q)=\chi_2(1/(q/2))$ and $\Ti_2(2/q)=\Ti_2(1/(q/2))$, so the
interest of \cref{thm:main} lies in odd $q$, where $2/q$ is not the reciprocal of an integer.
We have not read Hata's paper~\cite{Hata1990} ourselves. According to Krisman
\cite[\S1.2]{Krisman2026chi2}, whose reading we have not checked, Hata's linear independence
measures for shifted polylogarithms are stated at points $1/Q$ with $Q\in\Z$, so that for odd $q$
the relevant points $\pm4/q^2$ are not covered. The linear independence criterion of David,
Hirata-Kohno and Kawashima~\cite{DHK2022Acta} applies to the shifted dilogarithm
$\sum_{n\ge1}x^n/(n-\frac12)^2$ at rational points. As reported in \cite[\S1.2]{Krisman2026chi2},
Krisman evaluates it at $x=4/q^2$ and obtains the irrationality of $\chi_2(2/q)$ for
$q\ge 11\,795$ with the criterion as printed, and for $q\ge23\,587$ with a weaker form that he
believes the proof supports; these specialisations are his, not those of the authors of the
criterion. The point $x=-4/q^2$, which gives $\Ti_2(2/q)$, is also a rational point, so the same
criterion should give a threshold of the same order for $\Ti_2(2/q)$; we have not computed it and
have not found it stated. We did not search the classical literature on linear independence of
$1,\Li_2(x),\Li_2(-x)$ (Nikishin, Chudnovsky, Hata, Rhin--Viola, Marcovecchio and others)
ourselves; ``new'' below always means ``not found in the sources we read''.

The closest work is Krisman's preprint~\cite{Krisman2026chi2}, which claims proofs that
$\chi_2(1/q)\notin\Q$ for $q\ge3$, $q\ne4$, that $\Ti_2(1/q)\notin\Q$ for $q\ge5$, and that
$\chi_2(2/q)\notin\Q$ for odd $q\ge175$; it does not treat $\Ti_2(2/q)$ for odd $q$. That
preprint is labelled as primarily AI-generated and has not been refereed. We read version~2 and
checked against it the statement of each numbered result cited in this paper (Lemmas~2.1, 4.2, 5.1, 5.2,
6.4, 7.1, 11.1; Propositions~3.1, 11.2; Theorems~3.2, 4.5, 5.3, 8.1; Remark~6.7; \S1.2, \S12);
this was done by the AI system that revised this manuscript, not by a person, and it checks
statements, not proofs. We use none of its theorems as an input, and every estimate we need is
proved here. Its method, a Hankel-determinant construction
in the spirit of Zudilin's determinantal approach~\cite{Zudilin2017}, is the starting point of
ours: Krisman takes the Hankel determinants of the functional
$g\mapsto\sum_{t\ge1}w^t\,(g(u)\,u/D(u))'|_{u=t}$, where $D(u)=\prod_{k<h}(u+k+\frac12)$ has as
many half-integer poles as the size $h$ of the determinant, at $w=4/q^2$. In his \S12 he
states, numerically and assuming that $\frac{133}{60}$ is the true generic constant, that this
construction does not reach odd $q$ below about $87$: at the prime $2$ its leading coefficient
alone costs $\log2$ per $h^2$ \cite[Prop.~11.2]{Krisman2026chi2}, and with this cost and the
equilibrium constant the margin is positive only for odd $q\lesssim87$. We take this as a
reason to change the determinant rather than to sharpen the estimates for his.

\subsection{The constructions}\label{subsec:constructions}

Fix an odd integer $q\ge5$ (the hypothesis $q\ge5$ is a convenience; only $W<1$ is needed, that
is $q\ge3$, and the main theorem concerns $q\ge81$) and a sign $\sigma\in\{1,-1\}$, and put
\begin{equation}\label{eq:basic-notation}
  y=\frac q2,\qquad w=\frac{\sigma}{y^2}=\frac{4\sigma}{q^2},\qquad W=\abs w,\qquad
  \xi=\sum_{j\ge0}\frac{\sigma^jy^{-2j-1}}{(2j+1)^2},
\end{equation}
so that $\xi=\chi_2(2/q)$ for $\sigma=1$ and $\xi=\Ti_2(2/q)$ for $\sigma=-1$. For integers
$h\ge1$, $K\ge h$ and $A\ge0$ let
\[
  x_k=-k-\tfrac12\ (0\le k<K),\qquad D_K(u)=\prod_{k<K}(u-x_k),\qquad
  \varphi(u)=\frac{u\prod_{j=1}^{A}(u-j)^2}{D_K(u)}.
\]
In \cref{sec:pencils} we define an affine functional $\mu_X$ on the rational functions with
simple poles at the $x_k$ such that $\mu_\xi(R)=\sum_{m\ge1}w^mR'(m)$, and the polynomial
\[
  T(X)=\nrm\cdot\det\bigl[\mu_X(\varphi(u)\,u^{i+j})\bigr]_{0\le i,j<h}\in\Q[X],
  \qquad \nrm=\frac{(K!/A!^2)^h}{\prod_{i<h}i!^2},
\]
of degree at most $h$. We use two families of parameters $(h,K,A)$, indexed by $n\ge1$:
\begin{center}
\begin{tabular}{@{}lllll@{}}
\toprule
 & $(h,K,A)$ & $\kappa=K/h$ & $\alpha=A/h$ & used for\\
\midrule
general construction $\Gen$ & $(6n,\,7n,\,0)$ & $7/6$ & $0$ & odd $q\ge83$\\
band construction $\Band$ & $(40n,\,49n,\,2n)$ & $49/40$ & $1/20$ & $q=81$\\
\bottomrule
\end{tabular}
\end{center}
For $\Gen$ the only change with respect to Krisman's determinant is that there are $7n$ poles
for $6n$ rows. For $\Band$ the weight also vanishes to order two at $u=1,\dots,A$, which removes
the first $A$ terms of the series. We write $d=\kappa-2\alpha$, so $d=7/6$ for $\Gen$ and
$d=9/8$ for $\Band$.

\subsection{The argument in one page}\label{subsec:one-page}

For a prime $p$ and a nonzero $F\in\Q[X]$ let $\vG p(F)$ be the minimum of the $p$-adic
valuations of the coefficients of $F$. Let $\ell=\prod_pp^{-\vG p(T)}$, so that $\ell T\in\Z[X]$.
If $\xi=a/b$ were rational, then $b^h\ell T(a/b)$ would be an integer. The proof consists in
showing that this integer is nonzero for infinitely many $n$, while its absolute value is at
most $\exp\bigl(h^2(\bconst-\aconst+o(1))\bigr)$ with $\bconst<\aconst$
(\cref{prop:criterion}). This needs three estimates.

\begin{theorem}\label{thm:three-estimates}
Let $\mathcal C\in\{\Gen,\Band\}$, let $q\ge5$ be odd and $\sigma=\pm1$, and let
$T=T_n$ be the polynomial attached to the $n$-th member of $\mathcal C$, with $h=h_n$.
\begin{enumerate}[(i)]
\item \emph{Arithmetic.} Along the $n$ with $T_n\ne0$,
\[
  \liminf\frac{1}{h^2}\sum_p\vG p(T)\log p\ \ge\ \aconst_\sigma(q):=
  \Iprof_{\mathcal C}+c_{2,\mathcal C}\log2+\frac13\log q-\frac{(2-d)^2}4\log(q^2-4\sigma),
\]
with the constants $c_{2,\mathcal C}$ and $\Iprof_{\mathcal C}$ of \cref{tab:constants}.
\item \emph{Size.} For every comparison density $\rho$ and every $\lambda$ satisfying the tail
condition of \cref{thm:decay}, $\limsup h^{-2}\log\abs{T(\xi)}\le\bconst(\rho,\lambda)$, an
explicit number that depends on $q$ only through $W$.
\item \emph{Nonvanishing.} Let $r\in\Q$. If $\mathcal C=\Gen$ and $q\ge83$, or
$\mathcal C=\Band$ and $q=81$, then $T_n(r)\ne0$ for every sufficiently large prime $n$.
\end{enumerate}
\end{theorem}

\begin{proof}
Part~(i) is \cref{thm:arithmetic}, part~(ii) is \cref{thm:decay}, and part~(iii) is
\cref{thm:nonvanishing} together with \cref{prop:local-dets-general,prop:local-dets-band}.
\end{proof}

\begin{table}[ht]
\centering
\caption{The contributions to $\aconst_\sigma(q)$ and the comparison with $\bconst$, per
$h^2$. The constant $\aconst_-$ is the smaller of the two arithmetic constants
(\cref{rem:worse-sign}); $\bconst$ is the certified size exponent
(Appendix~\ref{app:decay-certificates}). Decimals are rounded in the safe direction; the first three rows are evaluated for $\sigma=-1$ and rounded to six places.}
\label{tab:constants}
\begin{tabular}{@{}lll@{}}
\toprule
 & $\Gen$, $q=83$ & $\Band$, $q=81$\\
\midrule
prime $2$ (\cref{thm:prime-two}) & \shortstack[l]{$-\frac{140}{99}\log2$\\$\approx-0.980208$} &
  \shortstack[l]{$-\frac{1847}{1100}\log2$\\$\approx-1.163857$}\\[4pt]
primes $p\mid q$ (\cref{prop:p-divides-q}) & \shortstack[l]{$+\frac13\log q$\\$\approx+1.472947$} &
  \shortstack[l]{$+\frac13\log q$\\$\approx+1.464816$}\\[4pt]
primes $p\mid q^2-4\sigma$ (\cref{prop:p-divides-Y}) &
  \shortstack[l]{$-\frac{25}{144}\log(q^2-4\sigma)$\\$\approx-1.534420$} &
  \shortstack[l]{$-\frac{49}{256}\log(q^2-4\sigma)$\\$\approx-1.682367$}\\[4pt]
ordinary primes, $\Iprof_{\mathcal C}$ (\cref{thm:ordinary}) & $\ge-2.113164137$ &
  $\ge-1.987089582$\\
\midrule
$\aconst_-(q)$ & $\ge-3.154846$ & $\ge-3.368498$\\
$\bconst$ & $\le-3.161708$ & $\le-3.374176$\\
margin $\aconst_--\bconst$ & $\ge0.006862$ & $\ge0.005678$\\
\bottomrule
\end{tabular}
\end{table}

The estimates are obtained as follows.
\begin{itemize}
\item \emph{Prime $2$} (\cref{sec:prime-two}). In the coordinate $z=2u$ the poles are odd and
the summation points even, and a Vandermonde-type minor whose nodes share a parity is divisible by
about $2^{t^2}$; the Cauchy--Binet formula then reduces $\vG2(T)$ to a quadratic minimisation in
the proportion $x$ of pole terms, with minimum at $x=29/33$ for $\Gen$ and $x=51/55$ for $\Band$.
\item \emph{Primes dividing $q(q^2-4\sigma)$} (\cref{sec:exceptional}). For $p\mid q$, reflecting
$m\mapsto-m$ makes the series $p$-adically convergent, and each level offers at most three terms
of valuation at least $2v_p(q)m$, which gives $\frac13v_p(q)h^2$; for $p\mid q^2-4\sigma$ the only
loss comes from the denominators $(1-w)^{b+1}$, and a convexity argument gives
$\frac{(2-d)^2}4v_p(q^2-4\sigma)h^2$.
\item \emph{Ordinary primes} (\cref{sec:ordinary}). For $p\asymp h$ two lower bounds, by Taylor
jets in each residue class modulo $p$ and by global partial fractions with Vandermonde gains, are
small minimisation problems over the residue classes; their maximum, integrated against $\log p$
by the prime number theorem, gives $\Iprof_{\mathcal C}$; the integral is evaluated exactly (Appendix~\ref{app:profile-certificate}).
\item \emph{Size} (\cref{sec:decay}). The determinant at $X=\xi$ is a sum over configurations of
$h$ points on the summation sites; Vandermonde bounds and a discrete energy inequality compare it
with an explicit step density, and the resulting number $\bconst$ is evaluated in ball arithmetic
(Appendix~\ref{app:decay-certificates}).
\item \emph{Nonvanishing} (\cref{sec:nonvanishing}). For $h=tp$ with $p$ prime, a basis adapted
to the residue classes modulo $p$ makes the scaled determinant block triangular modulo $p$, with
diagonal blocks equal, up to units, to three fixed local Hankel determinants, which are nonzero.
\item \emph{Assembly} (\cref{sec:proofs}). The margins of \cref{tab:constants} are positive; a
shift of $\lambda$ (\cref{prop:shift}) carries the case $q=83$ to every odd $q\ge83$. For each $q$
the inequality to be proved is $\bconst_q<\aconst_-(q)$, where $\bconst_q\le\bconst_{83}-\frac12\log\frac q{83}$
and $\aconst_-(q)\ge\aconst_-(83)-\frac1{72}\log\frac q{83}$; this is the whole uniformity
argument in $q$.
\end{itemize}

\subsection{What is new and what is inherited}\label{subsec:new}

The following ingredients come from Krisman's work~\cite{Krisman2026chi2}, in some cases in a
modified form: the functional and its pole series (\cref{lem:pole-series}); the irrationality
criterion (\cref{prop:criterion}); the reflection at primes dividing $q$ and the count of three
terms per level (\cref{prop:p-divides-q}); the degree count at primes dividing $q^2-4\sigma$
(\cref{prop:p-divides-Y}); the size bound by a discrete logarithmic gas compared with an explicit
step density, and its shift to smaller $\abs w$ (\cref{sec:decay}); and the strategy of proving
nonvanishing by reduction modulo a prime proportional to $h$ (\cref{sec:nonvanishing}). We reprove
all of them, because the constructions differ. What is new is the following.
\begin{enumerate}[(a)]
\item \emph{The ratio $K/h>1$.} More poles than rows lowers the degree of the polynomial part of
the weights. At primes dividing $q^2-4\sigma$ the loss drops from $\frac14\log(q^2-4\sigma)$ to
$\frac{25}{144}\log(q^2-4\sigma)$ per $h^2$, which at $q=83$ saves about $0.68$ and outweighs the
larger costs at $2$ and at the ordinary primes. With $K>h$ there is no Lagrange basis at the
poles, so the residue-class and nonvanishing arguments are set up with Taylor jets instead.
\item \emph{The prime $2$.} The parity argument gives the cost $\frac{140}{99}\log2$ for $\Gen$.
For Krisman's ratio $\kappa=1$ the limit of the bound is $\frac{13}{11}\log2\approx1.18\log2$,
which lies between his lower bound $\log2$ and his upper bound $\frac43\log2$
\cite[Lemma~11.1, Prop.~11.2]{Krisman2026chi2}. He also reports a measured loss of about
$1.18\log2$ per $h^2$ at $h=64$; we do not claim agreement with it (\cref{rem:two-comparison} and
Table~\ref{tab:two-adic-square} compare the finite values in two normalisations).
\item \emph{The ordinary primes.} The profile is valid for every $\kappa$ and with zeros in the
numerator. At $\kappa=1$ its integral equals Krisman's generic constant (\cref{rem:profile-k1}),
so it reorganises his class bound rather than improving it.
\item \emph{The zero band} (only for $q=81$). It costs in the size exponent and at the prime $2$,
but changes the balance between the remaining terms enough to reach $q=81$.
\item \emph{Nonvanishing.} The local determinants are $6\times6$ for $\Gen$, explicit
polynomials in $Q=\sigma q^2/4$ without roots of modulus $189$ or more, and $40\times40$ for
$\Band$, computed exactly at $q=81$.
\end{enumerate}

\subsection{Organisation and notation}\label{subsec:notation}

\Cref{sec:pencils} defines the functional, the determinants and the criterion.
\Cref{sec:prime-two,sec:exceptional,sec:ordinary} bound the denominators, \cref{sec:decay} the
size, and \cref{sec:nonvanishing} proves nonvanishing; \cref{sec:proofs} assembles the proof and
records the status of each step, and \cref{sec:limitations} explains why the method stops at $81$.
Appendix~\ref{app:certificates} describes the finite computations and
Appendix~\ref{app:consistency} compares the estimates with exact computations at small sizes.
\Cref{tab:notation} lists the notation used throughout.

\begin{table}[ht]
\centering
\caption{Notation.}\label{tab:notation}
\small
\begin{tabular}{@{}p{0.27\textwidth}p{0.66\textwidth}@{}}
\toprule
$q$, $\sigma$; $y$, $w$, $W$, $\xi$ & odd integer $\ge5$, sign $\pm1$; see~\eqref{eq:basic-notation}\\
$\tau_k$ & partial sum $\sum_{j\le k}\sigma^jy^{-2j-1}/(2j+1)^2$ of $\xi$, \eqref{eq:tau}\\
$h$, $K$, $A$; $\kappa$, $\alpha$, $d$ & size, number of poles, length of the zero band;
$K/h$, $A/h$, $\kappa-2\alpha$\\
$x_k$, $D_K$, $N_A$ & $-k-\frac12$; $\prod_{k<K}(u-x_k)$; $\prod_{j=1}^A(u-j)^2$\\
$\varphi$, $\Phi$ & $uN_A/D_K$; $(K!/A!^2)\,\varphi$, see~\eqref{eq:Phi}\\
$\Omega_b$, $\mom$, $\mu_X$ & $w^{\max(b,1)}(1-w)^{-b-1}$; rational moment sum; the functional
\eqref{eq:functional}\\
$\Delta$, $\nrm$, $T$ & Hankel determinant, normaliser, $T=\nrm\Delta$, see~\eqref{eq:Delta-T}\\
$\beta_i$, $\vvec(u)$, $\bvec(u)$ & $\binom ui$; the columns $(u^i)_{i<h}$ and $(\beta_i(u))_{i<h}$\\
$\vG p$ & Gauss valuation of a polynomial in $X$ over $\Q_p$\\
$\aconst_\sigma(q)$, $\bconst$, $\Iprof_{\mathcal C}$ & arithmetic constant, size exponent,
ordinary-prime constant (\cref{tab:constants})\\
$\psi(t)$, $\pos x$, $s_2(n)$ & $t\log\abs t-t$ with $\psi(0)=0$; $\max(x,0)$; binary digit
sum, so $v_2(n!)=n-s_2(n)$\\
\bottomrule
\end{tabular}
\end{table}

\emph{Conventions.} $p$ is always a prime and $v_p$ the $p$-adic valuation on $\Q_p$. The Gauss
valuation satisfies $\vG p(FG)=\vG p(F)+\vG p(G)$ and $\vG p(F+G)\ge\min(\vG p(F),\vG p(G))$, with
$\vG p(0)=+\infty$, and is invariant under $X\mapsto X+c$ for $c\in\Z_p$; hence the determinant of
a matrix over $\Q_p[X]$ has Gauss valuation at least the minimum over the terms of its Leibniz
expansion. For $u$ in any $\Q$-algebra, $\binom uk=u(u-1)\cdots(u-k+1)/k!$. The implied constants
in $O(\cdot)$ may depend on $q$ and on the construction, but not on $n$.
